\section{Introduction}
\label{sec:introduction}

The rapid integration of \gls{ai} into medicine faces a key challenge: clinicians have domain knowledge but lack technical skills for complex \gls{ml} pipelines, while data scientists can build advanced models but often lack the clinical context \cite{gisselbaek2025bridging}. This expertise gap hinders the adoption of \gls{ai} in clinical research.

\Gls{automl} platforms have emerged as promising solutions by automating complex tasks such as feature engineering, algorithm selection, and hyperparameter optimization, enabling users without extensive programming expertise to build predictive models \cite{karmaker2021automl}. However, current tools fail to bridge this gap for clinical researchers \cite{waring2020automated}. They treat datasets as generic matrices, ignoring structures inherent to clinical research, such as patient registries, observational studies, and clinical trials. They are not protocol-aware, making it difficult to distinguish between variables that define study populations (inclusion criteria) and those measuring outcomes (endpoints). Moreover, their black-box nature complicates interpretation by researchers. Therefore, incorporating domain knowledge into the modeling process is essential \cite{sirocchi2024medical}.

To address these limitations, we present \gls{hamelin}, a no-code, protocol-aware \gls{automl} platform grounded in \gls{hgml} principles, which emphasize integrating human expertise throughout the \gls{ml} pipeline. Its contribution lies in the design of a protocol-aware interface and the operationalization of \gls{hgml} principles in a clinical context, an area where, to the best of our knowledge, no prior system exists. \gls{hamelin} provides a collaborative environment where clinicians and data scientists work together using a shared language. By explicitly defining variables, distinguishing endpoints from exposures, and separating interventions from covariates, \gls{hamelin} facilitates the development of predictive models that are both clinically coherent and aligned with the structural requirements of research designs. This approach effectively bridges the gap between technical expertise and clinical domain.

% Revisar
The development of \gls{hamelin} followed a participatory design approach, involving \glspl{crp} from the Clinical Trials Agency of IDIVAL (Marqu\'es de Valdecilla University Hospital) throughout the design and refinement process. Four real clinical protocols served as use cases: EPA COVID-19 (a patient registry), UPS (an observational study on pressure injuries in neurosurgery), COLCHI-COVID (a clinical trial evaluating colchicine in COVID-19 patients), and VACAB-2022 (a clinical trial on bacterial vaccines for prosthesis infections). These protocols grounded the interface design in authentic clinical research scenarios and ensured that the system's terminology, workflow, and validation mechanisms align with actual clinical practice.

The remainder of this paper is organized as follows. Section~\ref{sec:background} provides background on clinical research methodologies, reviews related work in \gls{automl}, \gls{hcai}, \gls{hgml}, and introduces Ludwig as the underlying \gls{automl} engine. Section~\ref{sec:application} presents the participatory design process, the architecture, and the user interface \gls{hamelin}, detailing each component of the tab-based workflow. Section~\ref{subsec:hgml_alignment} evaluates the system's alignment with \gls{hgml} criteria, identifying strengths and limitations. Section~\ref{sec:discussion} discusses design decisions, trade-offs, and implications for future \gls{automl} systems. Finally, Section~\ref{sec:conclusions} concludes the paper and outlines directions for future work.
